% Options for packages loaded elsewhere
\PassOptionsToPackage{unicode}{hyperref}
\PassOptionsToPackage{hyphens}{url}
%
\documentclass[
  10pt,
]{article}
\usepackage{amsmath,amssymb}
\usepackage{iftex}
\ifPDFTeX
  \usepackage[T1]{fontenc}
  \usepackage[utf8]{inputenc}
  \usepackage{textcomp} % provide euro and other symbols
\else % if luatex or xetex
  \usepackage{unicode-math} % this also loads fontspec
  \defaultfontfeatures{Scale=MatchLowercase}
  \defaultfontfeatures[\rmfamily]{Ligatures=TeX,Scale=1}
\fi
\usepackage{lmodern}
\ifPDFTeX\else
  % xetex/luatex font selection
\fi
% Use upquote if available, for straight quotes in verbatim environments
\IfFileExists{upquote.sty}{\usepackage{upquote}}{}
\IfFileExists{microtype.sty}{% use microtype if available
  \usepackage[]{microtype}
  \UseMicrotypeSet[protrusion]{basicmath} % disable protrusion for tt fonts
}{}
\makeatletter
\@ifundefined{KOMAClassName}{% if non-KOMA class
  \IfFileExists{parskip.sty}{%
    \usepackage{parskip}
  }{% else
    \setlength{\parindent}{0pt}
    \setlength{\parskip}{6pt plus 2pt minus 1pt}}
}{% if KOMA class
  \KOMAoptions{parskip=half}}
\makeatother
\usepackage{xcolor}
\usepackage[margin=0.85in]{geometry}
\setlength{\emergencystretch}{3em} % prevent overfull lines
\providecommand{\tightlist}{%
  \setlength{\itemsep}{0pt}\setlength{\parskip}{0pt}}
\setcounter{secnumdepth}{-\maxdimen} % remove section numbering
\usepackage{amsmath,amssymb,amsthm,mathrsfs}
\ifLuaTeX
  \usepackage{selnolig}  % disable illegal ligatures
\fi
\usepackage{bookmark}
\IfFileExists{xurl.sty}{\usepackage{xurl}}{} % add URL line breaks if available
\urlstyle{same}
\hypersetup{
  pdftitle={Relative floors in big classes and the remaining finite-energy equation problem},
  pdfauthor={Research reconstruction for original P16.5},
  hidelinks,
  pdfcreator={LaTeX via pandoc}}

\title{Relative floors in big classes and the remaining finite-energy
equation problem}
\author{Research reconstruction for original P16.5}
\date{6 October 2026}

\begin{document}
\maketitle

\section{Abstract}\label{abstract}

We reconstruct an algebraic relative-floor argument for decreasing
closure of the genuine numerical plurisubharmonic cone in an arbitrary
big rational line-bundle class. The argument uses an actual
multiplier-ideal Fubini--Study net for each model obstacle, before any
closure theorem is invoked. Its compact comparison functions extend
upper semicontinuously across old poles. All tensor operations use
actual ample difference bundles. We give the directed-net construction,
an all-point endpoint proof, and the compactness argument, including the
primary ample proof of divisorial nonnegligibility. We then reconstruct
a target-independent section-volume energy and its entire-cone outer
derivative construction, with raw mass equal to volume. The final
sections identify the precise further assertions required for the
equation, audit the historical proposed passages, and distinguish the
proved zero numerical class from other zero-volume classes. This paper
does not claim a completed original P16.5 theorem. All conclusions are
research drafts, and the dependency qualifications in the source ledger
are part of their statements.

\section{1. Fixed carrier, actual bundles, and available
premises}\label{fixed-carrier-actual-bundles-and-available-premises}

Let \(k\) be any algebraically closed field of characteristic zero with
the trivial absolute value. Let \(X\) be smooth, integral, projective,
and of dimension \(n\geq1\). Write \(K=X^{\mathrm{an}}\), with its
compact Hausdorff Berkovich topology. Fix an \textbf{actual} big
rational line bundle \(L\), an actual ample rational line bundle \(A\),
and
\[\theta=c_1(L),\qquad a=c_1(A),\qquad V=\operatorname{vol}_X(L)>0.\] No
hypothesis concerning cardinality, finite generation, nefness, torus
invariance, support, or a chosen skeleton is imposed.

For an actual rational bundle \(B\), a generically finite Fubini--Study
function is \[\phi=\frac1m\max_j\{\log|s_j|+\lambda_j\},\qquad
s_j\in H^0(X,mB)\setminus\{0\},\quad \lambda_j\in\mathbb R.\] Here
\(mB\) is an honest bundle and the norms use its canonical trivial
metric. Such a function is continuous into \(\mathbb R\cup\{-\infty\}\);
its generic value is finite and is its supremum. The cone
\(\operatorname{PSH}(\theta)\) consists of all all-point decreasing
limits of these functions on actual bundles \(B_i\) satisfying
\(c_1(B_i)\to\theta\), with finite limiting generic value. The topology
on this cone is pointwise convergence at rational divisorial valuations
and the generic trivial valuation \(v_0\). We do not define the cone or
its finite-energy part as a completion.

The exact supplied G0 body is used only for actual ample bundles:
fixed-bundle finite Fubini--Study regularization, normalized
compactness, finite divisorial values, generic normalization, and
all-point order determination. Its equation is never applied to \(L\) or
another nonample bundle. All other mathematical inputs are external
algebraic or actual-ample inputs and are recorded separately. In
particular, general big compactness, the abstract twice-closure lemma,
and the Addendum's general envelope consequence are not premises.

We use the ordinary algebraic volume theorem on smooth projective
varieties: true full-section asymptotics for big rational bundles,
birational invariance and monotonicity, continuity of volume, and
\[d\operatorname{vol}_{Y,\alpha}(H)
=d\langle\alpha^{d-1}\rangle\cdot H,
\quad d=\dim Y,\] on the big cone. Positive products are homogeneous and
monotone in the pseudoeffective order; a positive Weil curve-class
difference has nonnegative pairing with a nef Cartier class. These are
not non-Archimedean equation grants. Their precise proof units and
remaining background dependencies are recorded in the ledger. A theorem
asserted below relative to these algebraic inputs is a mathematical
conditional deduction, not a computational certification.

\section{2. Elementary operations on actual
bundles}\label{elementary-operations-on-actual-bundles}

If \(C\) is an actual ample rational bundle, choose a divisible globally
generated power. The maximum of the logarithms of a generating system is
identically zero: at the center of any analytic point some generator is
a unit. Tensoring an FS system on \(B\) with that generating system
represents exactly the same scalar function on \(B+C\).

If \(c_1(B_i)\to\theta\) and \(\Omega-L\) is ample, then \(\Omega-B_i\)
is an actual ample rational bundle for all sufficiently late \(i\). Thus
the tail of any defining FS net may be represented in the same actual
ample \(\Omega\). Appending a generating system of \(m\Omega\) at weight
\(-mk\) represents \(\max(\phi_i,-k)\) in that bundle. The product net
\((i,k)\) decreases to the same endpoint, including at poles. Therefore
every genuine \(\theta\) function embeds as the same original-carrier
function into a fixed ample cone. This does not assert that its
\(\theta\) energy is its ample energy.

Finite maxima are genuine without general closure. For defining nets for
\(u,v\), choose actual common targets \(B_{i,j}\) whose differences from
the two incoming bundles are ample and whose classes tend to \(\theta\).
Tensor both systems into that \textbf{one} target, and take their union.
The represented scalar function is \(\max(\phi_i,\psi_j)\), independent
of these choices, so it decreases on the product index.

Real convex combinations have one defining net as well. On a tail
\(\phi_i,\psi_j\leq C\), choose positive rational \(r_\ell\uparrow r\)
and \(s_\ell\uparrow1-r\). The functions
\[r_\ell(\phi_i-C)+s_\ell(\psi_j-C)+C\] decrease on the product index,
are actual FS functions after clearing rational tensor powers, and have
limiting numerical class \(\theta\). Their endpoint is \(ru+(1-r)v\). No
subtraction of two poles occurs. Constants and positive rational scaling
are immediate at finite stages.

For normalized functions the fixed ample embedding gives, at every fixed
rational divisorial \(z\), \[-T_\Omega(z)\leq u(z)\leq0.\] The constant
is finite. At an FS stage the lower bound follows from the linear-growth
estimate \(v(s)\leq mT_\Omega(v)\) for every nonzero section of
\(m\Omega\), and then passes through the given defining net. This
section estimate, or the exact ample G0 bound, suffices; no nonample
compactness is involved.

\section{3. One model-obstacle multiplier-ideal
net}\label{one-model-obstacle-multiplier-ideal-net}

This section reconstructs the independent algebraic construction needed
to produce references. Let \(B\) be an actual big rational bundle and
\(f\) a rational model function. Represent \(f\) by a test configuration
\((\mathcal X,\mathcal B)\) with generic bundle \(B\); resolve the total
space and its central pair without altering the punctured product. Clear
denominators. Over \(\mathbb A^1\), let \(\mathfrak a_m\) be the
\textbf{relative complete-system} base ideal of \(m\mathcal B\). These
are ideals on that smooth total space, not base ideals on an arbitrarily
substituted carrier.

Normality gives the section criterion: for \(s\in H^0(X,mB)\), the
monomial \(s t^{-\lambda}\) extends over the test configuration if and
only if \[\log|s|+\lambda\leq m f\quad\hbox{at every point of }K.\]
Indeed the regularity inequality at each central prime is precisely this
inequality at its normalized Gauss valuation. Regularity away from the
central fiber is already known. Nonnegative prime orders imply
regularity on the normal total space. Conversely a regular local section
has nonnegative value at every centered semivaluation. This proves the
criterion also at proper-support points, not only at valuations.

It follows that
\[f_m=f+m^{-1}\phi_{\mathfrak a_m}\in\operatorname{FS}(B)\] increases in
the divisibility order and has pointwise supremum \(Q_B(f)\), the
supremum of actual FS minorants in \(B\). Multiplication gives
\(\mathfrak a_m^\ell\subset\mathfrak a_{m\ell}\). For the opposite
inequality, any finite weighted section system below \(f\) consists of
extending monomials in some degree, hence is below the corresponding
\(f_m\). Real constants can be approximated by rational constants from
below.

Set \(\mathfrak b_m=\mathcal J(\mathfrak a_\bullet^m)\). The
multiplier-ideal inputs are \[\mathfrak a_m\subset\mathfrak b_m,
\qquad \mathfrak b_{m\ell}\subset\mathfrak b_m^\ell,
\qquad (m\mathcal B+\mathcal C)\otimes\mathfrak b_m
\text{ globally generated for every }m,\] for one actual equivariant
ample \(\mathcal C\). Their domain is smooth characteristic-zero total
spaces and relative complete systems. They are not consequences of a
closure lemma.

Here is the algebraic mechanism and its explicit background boundary.
Asymptotic multiplier ideals are the maximal stabilized ideals
\(\mathcal J((m/q)\mathfrak a_q)\) in a divisible tail. For a common
sufficiently divisible \(q\), ordinary ideal subadditivity proves the
second inclusion. The first uses
\(\mathfrak a_m^q\subset\mathfrak a_{mq}\) and
\(\mathfrak a_m\subset\mathcal J(\mathfrak a_m)\) on a smooth variety.
The diagonal/product proof of subadditivity is algebraic, after the
restriction theorem for multiplier ideals. Each individual comparison
uses finitely many ideals and resolutions, so finite-data descent to a
characteristic-zero subfield and flat extension apply; an entire
infinite sequence need not descend to one countable field.

For uniform relative generation choose a relatively very ample
equivariant \(H\) on the smooth total space of dimension \(d=n+1\),
sufficiently positive that \[\mathcal C=K_{\mathcal X}+(d+1)H\] is
ample. Resolve a sufficiently divisible base ideal computing
\(\mathfrak b_m\). Its moving part is relatively generated. Relative
Kawamata--Viehweg/local vanishing, followed by the multiplier-ideal
projection formula, gives the relative cohomology vanishings for
\[\mathcal O(m\mathcal B+\mathcal C-iH)\otimes\mathfrak b_m,
\qquad 1\leq i\leq d.\] Relative Castelnuovo--Mumford regularity then
gives generation. Since the base is affine, this is global generation
over \(\mathbb A^1\). The bundle is chosen once, independently of \(m\).
Characteristic-zero relative vanishing, local vanishing, and the
restriction theorem remain ordinary external algebraic inputs: the
retrieved DEL proof is complex and absolute; it must not be labelled a
direct proof of the relative arbitrary-field statement. This exact
distinction is an outstanding primary-source obligation in the final
ledger. The conditional construction below retains this input
explicitly.

Shift the central linearization of \(\mathcal C\) so that its model
function \(\phi_{\mathcal C}\geq0\). The genuine actual FS stages
\[\Psi_m=f+m^{-1}\phi_{\mathfrak b_m}
+m^{-1}\phi_{\mathcal C}
\in\operatorname{FS}(B+C/m)\tag{3.1}\] decrease in the divisibility
order. They satisfy \[Q_B(f)\leq\Psi_m\leq f+m^{-1}\phi_{\mathcal C}.\]
A shifted nonzero \(B\) section below the bounded obstacle prevents
collapse of their generic values. Their all-point limit \(\Psi\)
therefore belongs to \(\operatorname{PSH}(c_1(B))\) \textbf{directly by
its one defining net}. No decreasing closure has yet been used.

To identify the limit with the genuine envelope, let \(\tau\leq f\) be
genuine \(B\)-psh. Choose an actual decomposition \(B=C_0+E\) with
\(C_0\) ample and \(E\) effective. Its existing defining net can, by
actual ample difference tensoring, be represented without changing
values in \(B+\epsilon_jC_0\), where \(\epsilon_j\to0\). Compactness
against the finite continuous \(f+\epsilon\) makes a sufficiently late
tail at most \(f+\epsilon\). Then
\[\frac{\tau_j+\epsilon_j\log|s_E|}{1+\epsilon_j}
\in\operatorname{FS}(B).\] The actual bundle identity is
\((B+\epsilon_jC_0)+\epsilon_jE=(1+\epsilon_j)B\). Since \(f\) is
bounded and \(\log|s_E|\leq0\), this function is at most
\(Q_B(f)+\epsilon+O(\epsilon_j)\). On \((X-E)^{\mathrm{an}}\) the
divisor logarithm is finite. Passing to the limit proves
\(\tau\leq Q_B(f)\leq\Psi\) at every divisorial point. The common actual
ample embedding and G0 order determination give \(\tau\leq\Psi\)
everywhere. Thus \(P_B(f)=\Psi\) on \(K\).

In particular \(V_B=P_B(0)\) exists, has generic value zero, and
dominates every normalized genuine \(B\) potential. It is produced by
(3.1), not by general closure or compactness. Normalizing any one stage
\(\psi\) of that particular net gives
\[\psi\leq0,\qquad V_B\leq\psi+c_\psi\tag{3.2}\] with a finite constant
\(c_\psi\).

\section{4. The comparison function at old
poles}\label{the-comparison-function-at-old-poles}

\textbf{Lemma 4.1.} Suppose \(\psi\leq0\) is continuous into
\(\mathbb R\cup\{-\infty\}\), \(0<\delta<1\), and an old FS function
satisfies \[F\geq(1-\delta)\psi-k.\] Let \(h\) be usc and satisfy
\(h\leq\psi+C\) for a finite constant \(C\). Define \[D_F(h)(x)=
\begin{cases}h(x)-F(x),&F(x)>-\infty,\\
-\infty,&F(x)=-\infty.
\end{cases}\] Then this is usc. The assignment at old poles is a
comparison extension; it is not a definition of \((-\infty)-(-\infty)\).

\emph{Proof.} Off the old poles the denominator is finite continuous
locally and \[D_F(h)\leq\delta\psi+C+k.\] At a pole of \(F\), the floor
inequality forces \(\psi=-\infty\), and the upper bound forces
\(h=-\infty\). Continuity of \(\psi\) into the extended half-line makes
the displayed upper bound tend to \(-\infty\) in every neighborhood
filter at that point. The assigned value is therefore usc. If \(h_i\)
decreases, the extensions decrease pointwise. If their limiting
comparison is strictly negative everywhere, their nested closed
nonnegative superlevel sets have empty intersection. Compactness and
directedness make one such set empty and every later set empty. This
works for arbitrary directed nets. \(\square\)

The uniform upper bound is essential. Continuity into the extended
half-line, constant shifts, and finite generic value alone do not supply
it. The abstract twice-decreasing-closure counterexample in the supplied
review remains valid and is reproduced in Section 12 below.

\section{5. The actual directed FS
construction}\label{the-actual-directed-fs-construction}

\textbf{Theorem 5.1 (relative-floor closure, with the explicit algebraic
input of Section 3).} If \(u_i\in\operatorname{PSH}(\theta)\) decreases
and \(u=\inf_i u_i\) has finite generic value, then \(u\) has one
all-point decreasing FS defining net with classes tending to \(\theta\).

\emph{Proof.} Restrict to a tail with \(u_i\leq M\), a finite common
bound. For each positive rational \(\eta\) use the actual target bundle
\(T_\eta=L+\eta A\). Choose \(0<\epsilon<\eta/4\) and the particular
reference \(V_B\) for \(B=L+\epsilon A\). Choose a normalized stage
\(\psi\) of (3.1), with actual bundle \(R\) and class \(\beta=c_1(R)\),
sufficiently late that \(T_\eta-R\) is ample. Choose rational
\(0<\delta<1\) so small that
\[T_\eta-(1-\delta)R\quad\hbox{is an actual ample rational bundle.}\]
This condition is checked in the open ample cone: before choosing
\(\delta\), \((\eta-\epsilon)a-c_1(C)/m\) is ample for large \(m\). Its
sufficiently small perturbation by \(\delta c_1(R)\) stays ample.
Equality of numerical classes is never substituted for an actual tensor
identification.

A special stage is an FS function \(F\) on this \textbf{specified
actual} \(T_\eta\) with some data as above, some finite \(k\), and a
positive rational guard \(\kappa\), satisfying \[F\geq u+\kappa,
\qquad F\geq(1-\delta)\psi-k.\tag{5.1}\] Such functions exist: take a
late FS upper stage \(h\) of some \(u_i\), whose actual bundle
difference from \(T_\eta\) is ample, tensor it into \(T_\eta\), tensor
the rationally scaled floor into the same bundle, and take
\[F=\max(h+\kappa,(1-\delta)\psi-k).\] Clearing rational tensor powers
and taking the union of the two section lists gives one actual FS
function.

Every \(u_i\) embeds into \(B\) and hence \(u_i\leq V_B+M\). A
sufficiently late tail of the defining FS stages of any fixed \(u_i\)
embeds into \(B\) too; its generic values are at most \(M+1\). Thus both
kinds of incoming functions satisfy
\[h\leq V_B+M+1\leq\psi+c_\psi+M+1.\tag{5.2}\] The constant is uniform
for a fixed old reference. This proves the algebraic hypothesis of Lemma
4.1 for \textbf{all late incoming stages}, including at the old
reference poles.

Take finitely many old special stages \(F_\ell\) with data
\((\eta_\ell,B_\ell,\psi_\ell,\delta_\ell,k_\ell,\kappa_\ell)\). Choose
a new \(\eta>0\) as small as desired and a new reference stage
\(\psi\leq0\) sufficiently near \(\theta\) that its actual bundle \(R\)
satisfies \(B_\ell-R\) ample for every \(\ell\). It is a normalized
genuine \(B_\ell\) potential, so
\[\psi\leq V_{B_\ell}\leq\psi_\ell+c_\ell.\] Choose
\(\delta<\min_\ell\delta_\ell\) sufficiently small for the new target
amplitude condition. For \(b=(1-\delta)\psi-k\),
\[b-F_\ell\leq(\delta_\ell-\delta)\psi_\ell
+(1-\delta)c_\ell+k_\ell-k\] where \(F_\ell\) is finite. The first term
is nonpositive. Increasing \(k\) makes \(b\leq F_\ell\) simultaneously.
At an old pole the preceding literal inequality between the references
already gives \(b=-\infty\), so the order holds there too.

Choose \(0<\kappa<\frac14\min_\ell\kappa_\ell\). Apply Lemma 4.1 to
\(u_i+2\kappa\) and each old stage. At finite old values the limiting
comparison is at most \(2\kappa-\kappa_\ell<0\); its extension at old
poles is \(-\infty\). A single late \(i\) satisfies
\(u_i+2\kappa<F_\ell\) off all old poles. Now apply the same compact
argument to a sufficiently late defining FS tail of this fixed \(u_i\),
using (5.2) separately for each old reference. Choose its stage
\(h\geq u_i\geq u\) with
\[h+\kappa\leq F_\ell\quad\hbox{for every }\ell,\] and with actual
bundle close enough to \(L\) to tensor into the new \(T_\eta\). Then
\[F=\max(h+\kappa,b)\leq F_\ell\quad\hbox{for every }\ell,
\qquad F\geq u+\kappa.\] Thus the special stages are downward directed,
while their positive \(\eta\) coordinates can be made arbitrarily small
after any finite list of constraints. Ordering the data by scalar
decrease and \(\eta\) decrease yields one directed net. Different data
representing equal functions cause only a preorder; quotienting equal
data is unnecessary.

Its infimum is exactly \(u\) at every point. Fix \(x\) and a real
threshold \(t>u(x)\). First choose \(u_i(x)<t\) with some positive
margin, and \textbf{choose \(\eta\) before taking a sufficiently late
defining stage \(h\)} so that its actual bundle fits \(T_\eta\). Then
take \(h(x)\) close enough to \(u_i(x)\) and the guard small enough that
\(h(x)+\kappa<t\). If \(\psi(x)\) is finite, increase \(k\) to put the
floor below that value; if it is a pole, the floor is already
\(-\infty\). This gives \(F(x)<t\). If \(u(x)=-\infty\), use any finite
\(t\) and the same argument. To impose old constraints as well,
directedness gives a later stage below this one and the old stages. The
limit is therefore all-point \(u\), and its classes
\(c_1(T_\eta)=\theta+\eta a\) tend to \(\theta\). Its generic value is
the stipulated finite one. \(\square\)

The positive guard applies to the incoming \textbf{upper approximation}.
Multiplying that incoming stage by \(1+\delta\) would not preserve the
inequality above the desired endpoint and is not used. The factor
\(1-\delta\) weakens only the old-pole reference floor.

\section{6. Envelopes and compactness, without a big compactness
import}\label{envelopes-and-compactness-without-a-big-compactness-import}

Let \(f\) be usc and bounded above, and suppose it has a genuine
\(\theta\) minorant with finite generic value. Continuous real majorants
decrease to \(f\) pointwise on a compact Hausdorff space. Uniform
rational model approximation from above, and finite minima, give a
directed net of rational model majorants with the same infimum. Their
genuine envelopes were constructed in Section 3. They decrease, and
their generic values stay bounded below by the minorant's finite value.
Theorem 5.1 gives a genuine all-point limit \(v\leq f\). Every
competitor for \(f\) is below each of these envelopes, while \(v\) is
itself a competitor. Hence \(v=P_\theta(f)\).

For a bounded-above family of genuine functions, apply this construction
to its usc supremum. The envelope dominates the raw supremum and, being
usc, also its usc regularization. The obstacle inequality gives
equality. Thus the envelope property follows from the independently
constructed model envelopes and the new relative-floor closure.

We supply the actual-ample nonnegligibility proof needed for
compactness. Fix an ample \(\Omega\) and a rational divisorial \(z\).
Choose a flag ideal whose normalized blowup has \(z\) among its finite
Rees set \(S\). For a normalized \(\Omega\) potential, values at the
finitely many points of \(S\) have a uniform finite spread bound, by
Section 2.

The algebraic separator construction proceeds as follows. On the normal
flag blowup, twisting its relatively ample antieffective divisor by a
sufficiently large multiple makes both it and its twist by one central
prime relatively generated. Their pushforward ideals have a strict
inclusion. After twisting by a large actual multiple of \(\Omega\),
select sections whose normalized values are strictly larger than the
flag function at one selected Rees point and no larger at the others.
Appending a finite constant floor makes their maxima finite FS
functions. Subtract the value gaps, add prescribed rational constants,
and take their maximum. Multiplying the resulting difference from the
flag function by a sufficiently large positive integer supplies a finite
continuous model function \(q\) with
\[\sup_K(w-q)=(w-q)(z)\quad\hbox{for every }w\in\operatorname{PSH}(\Omega).\tag{6.1}\]
The finite-vertex maximum property for FS functions follows directly
from the regularity criterion on the normal flag blowup. For an
arbitrary ample \(w\), take its G0 fixed-ample decreasing finite FS net
and pass the finite maximum over \(S\) through the limit. This is the
full step that the primary proof calls ample regularization. It uses G0
in its retained ample domain and never invokes big decreasing closure.
The primary complete construction is BJ Lemmas 2.28--2.29, Lemma 4.26,
Lemma 5.7, and Theorem 5.6, at the recorded literal coordinates.

For a bounded-above family \(w_i\) in this ample cone, (6.1) gives
\[\sup_i w_i\leq\bigl(\sup_i w_i(z)\bigr)-q(z)+q.\] Since \(q\) is
finite continuous, its usc regularization satisfies the same inequality.
Evaluating at \(z\) proves
\[\left(\sup_i w_i\right)^*(z)=\sup_i w_i(z).\tag{6.2}\] Thus rational
divisorial points are nonnegligible in the actual ample cone. This is a
mathematical proof through the primary separator construction, not an
assumption that compactness implies nonnegligibility.

Now embed the normalized \(\theta\) cone in the one normalized actual
ample \(\Omega\) cone. Let a net of its elements converge divisorially
in that ample compact space. Its regularized tail suprema are genuine
\(\theta\) potentials by the envelope property proved above. They
decrease; (6.2), applied in \(\Omega\), identifies their divisorial
values with the suprema of the tail coordinates. Their generic values
remain zero. Theorem 5.1 makes their all-point infimum genuine
\(\theta\); its coordinates are the prescribed limit. Ample all-point
determination identifies that original function with the ample limit.
Hence the normalized \(\theta\) cone is closed in the compact ample
cone, and is compact. No Addendum Corollary 2.2 or BJ general Theorem
5.11 is imported for this conclusion.

\section{7. Actual section lattices and volume
mass}\label{actual-section-lattices-and-volume-mass}

Let \(\rho:Y\to X\times\mathbb P^1\) be a smooth equivariant
modification, unchanged off zero, and let
\[F=\sum_i b_iE_i=\rho^*\{0\},\quad D=\sum_i d_iE_i,
\quad \alpha_D(c)=\rho^*p_X^*L+D+cF.\] All \(b_i>0\) are retained. A
vertical model function \(f_D\) has value \(d_i/b_i\) at
\(x_i=b_i^{-1}\operatorname{ord}_{E_i}|_{k(X)}\). Its scalar constant
shift is represented by \(D+cF\).

For divisible \(m\), set \(R_m=H^0(X,mL)\) and \(r_m=\dim R_m\). The
actual module \[M_m=H^0(Y|_{\mathbb A^1},m(\rho^*L+D))
\subset R_m\otimes k[t,t^{-1}]\] is a graded full \(k[t]\) lattice. If
\(aF\leq D\leq bF\) coefficientwise, it lies between \(t^{-ma}R_m[t]\)
and \(t^{-mb}R_m[t]\). An adapted basis gives intrinsic weights
\(w_{m,j}\in\mathbb Z\), \(ma\leq w_{m,j}\leq mb\). Glue to the fixed
product lattice at infinity. Its pushforward is a vector bundle on
\(\mathbb P^1\) and splits as \(\bigoplus_j\mathcal O(w_{m,j})\). The
torsion-free direct image is locally free because the base is a
nonsingular curve. Consequently
\[h^0(Y,m\alpha_D(c))=\sum_j\max(w_{m,j}+mc+1,0).\tag{7.1}\] For
\(c>-a\) this is exactly \((mc+1)r_m+\sum_jw_{m,j}\). The full-section
limit \(n!r_m/m^n\to V\) and the true volume limits on \(Y\) imply, for
every sufficiently large \(c\),
\[\operatorname{vol}_Y(\alpha_D(c))=(n+1)cV+(n+1)\mathscr E(D).\tag{7.2}\]
The intercept \(\mathscr E(D)\) is independent of \(c\) and of a higher
model, by birational volume invariance. It is monotone, concave, and
satisfies \[\mathscr E(D+sF)=\mathscr E(D)+sV,
\qquad |\mathscr E(D)-\mathscr E(G)|\leq V\|f_D-f_G\|_\infty.\]
Concavity follows by applying concavity of the \(1/(n+1)\) power of
volume to (7.2), expanding at large \(c\), and letting \(c\to\infty\).

For large \(c\) define \(q_i=\langle\alpha_D(c)^n\rangle\cdot E_i\). It
is nonnegative: adding the effective \(E_i\) increases volume, and its
right derivative equals \((n+1)q_i\). Uniformly for a small change of
\(D\) in the \(E_i\) direction, (7.2) holds on one common half-line.
Differentiating there proves that each \(q_i\) is independent of large
\(c\), not merely that its sum is independent. Also \[\sum_i b_iq_i=V.\]
Thus \[\mu_D=\sum_i b_iq_i\,\delta_{x_i}\tag{7.3}\] is an actual
positive original-carrier atomic measure of raw mass
\(\operatorname{vol}(L)\).

On a higher model, effective exceptional divisors \(G\) satisfy
\(f_*\mathcal O(kG)=\mathcal O\) by normality of the base. Hence adding
them to a pulled big bundle leaves section counts and volume unchanged.
Its volume differential has zero exceptional coefficient. Strict
transforms have the old coefficient by the projection formula. This
proves invariance of (7.3) under model pullback, including all central
multiplicities.

For every vertical \(H\) on a common model,
\[\left.\frac{d}{dt}\right|_{0}\mathscr E(D+tH)
=\int f_H\,d\mu_D.\tag{7.4}\] This construction precedes any target
choice. It does not yet identify the measure with the finite-energy
Monge--Ampère measure of an arbitrary endpoint.

\section{8. Continuous obstacles and entire genuine
energy}\label{continuous-obstacles-and-entire-genuine-energy}

For a continuous obstacle \(f\) and a nonzero actual section
\(s\in R_m\), set \[W_m(f;s)=\inf_{x\in K}\{mf(x)-\log|s|(x)\}.\] The
bracket is \(+\infty\) at a section pole. It has a finite minimum
elsewhere. Its superlevel sets, with zero adjoined, are vector subspaces
by the non-Archimedean triangle inequality. The decreasing filtration
has at most \(r_m\) nontrivial jumps, and closed endpoint conditions
follow from its intersection definition. Its intrinsic weights are
therefore well defined with multiplicities. Their floored mean
\[e_m(f)=\frac{n!}{m^{n+1}}\sum_j\lfloor W_m(f;s_{m,j})\rfloor\] does
not depend on an adapted basis.

For model \(f_D\), the section criterion identifies these floors with
the lattice weights in Section 7, so \(e_m(f_D)\to\mathscr E(D)\). If
\(\|f-g\|\leq\epsilon\), filtration comparison gives
\[|e_m(f)-e_m(g)|\leq \frac{n!r_m}{m^n}(\epsilon+1/m).\] Uniform model
density thus proves existence of the limit \(\mathscr E(f)\) for every
\(f\in C(K)\) and identifies it with the unique \(V\)-Lipschitz
extension of the model energy.

Every section term \(m^{-1}\log|s|+\lambda/m\) is a genuine \(\theta\)
minorant. Consequently \[W_m(f;s)\geq\lambda
\quad\Longleftrightarrow\quad
m^{-1}\log|s|+\lambda/m\leq P_\theta(f).\tag{8.1}\] This is all-point,
including section zeros. It proves that \(\mathscr E(f)\) depends only
on \(P_\theta(f)\), and is ordered by that envelope. Uniform original
continuous obstacle approximation never assigns finite energy to an
uncapped fixed section system.

For an actual genuine function define
\[E_\theta(u)=\inf\{\mathscr E(f):f\in C(K),\ f\geq u\},
\qquad \mathcal E^1(\theta)=\{u\in\operatorname{PSH}(\theta):E_\theta(u)>-\infty\}.
\tag{8.2}\] It agrees with \(\mathscr E(f)\) on \(P_\theta(f)\), by
envelope order. It is monotone and shifts by \(cV\). Its reference is
the genuine \(V_\theta=P_\theta(0)\) from Section 3, with
\(E_\theta(V_\theta)=0\).

If \(u_i\downarrow u\) genuinely, the energy decreases to
\(E_\theta(u)\). Indeed, for continuous \(f\geq u\) and \(\epsilon>0\),
the nested compact sets \(\{u_i\geq f+\epsilon\}\) have empty
intersection. Eventually \(u_i\leq f+\epsilon\), whence
\(\inf E_\theta(u_i)\leq\mathscr E(f)+\epsilon V\). Infimize over \(f\)
and let \(\epsilon\downarrow0\). This also covers energy \(-\infty\).

Normalized \(u\) satisfies \(u\leq V_\theta\). The relative floors
\[u_k=\max(u,V_\theta-k)\] are genuine, decrease all-point to \(u\), and
obey \(-kV\leq E_\theta(u_k)\leq0\). Thus (8.2) is equivalently finite
limiting relative-floor cost. On the reference pole set both branches
are \(-\infty\); no subtraction is defined there. Finite energy does
\textbf{not} mean bounded difference from \(V_\theta\).

At a model base, concavity and (7.4) give a global supporting plane on
continuous obstacles. Passing continuous majorants down to a genuine
\(u\) is legitimate at the finitely many divisorial atoms, where all
values are finite. This yields \[E_\theta(u)\leq A_D(u):=
\mathscr E(f_D)+\int(u-f_D)\,d\mu_D,
\qquad E_\theta(u)=\inf_D A_D(u).\tag{8.3}\] For the equality, a model
majorant of \(u\) gives \(A_D(u)\leq\mathscr E(f_D)\), and such
majorants are cofinal. The index \(D\) in (8.3) is independent of \(u\).
Each \(A_D\) is a finite-coordinate affine functional, so the genuine
energy is concave and weakly upper semicontinuous. By Section 6,
normalized sets \(E_\theta(u)\geq-B\) are compact and consist of actual
finite-energy functions.

\section{9. Truncated volumes and a target-independent entire outer
construction}\label{truncated-volumes-and-a-target-independent-entire-outer-construction}

This section records the entire operator deduction conditional on the
exact algebraic inputs and genuine foundations above. It is more than an
inverse defined by targets, but is not yet the full equation theorem.

For a nonpositive model \(f_D\) of cost \(-\mathscr E(D)\leq B\), put
\[g_D(c)=\frac1{n+1}\operatorname{vol}_Y(\alpha_D(c)).\] Equation (7.1),
including zero volume when a class is not big, gives
\[g_D(c)=\lim_m\frac{n!}{m^{n+1}}
\sum_j(w_{m,j}+mc)_+.\] All \(w_{m,j}\leq0\). Hence \(g_D\) is convex,
\(V\)-Lipschitz, \(g_D(0)=0\), and \(g_D(c)\geq cV-B\). For \(c>B/V\),
\(\alpha_D(c)\) is big and
\[\mu_{D,c}=\sum_i b_i\langle\alpha_D(c)^n\rangle\cdot E_i\,\delta_{x_i}\]
is positive, of mass \(g_D'(c)\in[V-B/c,V]\).

The decisive estimate is coefficientwise: \[0\leq\mu_{D,c}\leq\mu_D,
\qquad\|\mu_D-\mu_{D,c}\|_{\mathrm{TV}}\leq B/c.\tag{9.1}\] To prove it,
write \[\mathscr E(D)+cV-g_D(c)
=\lim_m\frac{n!}{m^{n+1}}\sum_j\min(w_{m,j}+mc,0).\] Adding an effective
vertical prime orders every section flag and makes this expression
increase. Differentiate in that prime direction, using the eventual
intercept derivative and the fixed-\(c\) big volume derivative. The
result is exactly the nonnegative difference between the two
coefficients. This avoids pairing a positive curve-class order with a
nonnef central divisor.

Fix a model direction \(h=f_H\) on one smooth model \(Z\) and fixed
\(c>B/V\). On higher common models let \(\gamma=\rho_Z^*L+cF\),
\(\alpha=\alpha_D(c)\), and \(T\) be a fixed ample Cartier class on
\(Z\). Then \(\alpha\leq\gamma\) and
\(\operatorname{vol}(\alpha)\geq(n+1)(cV-B)\). While \(\alpha-rT\) is
big, monotonicity paired with the nef pulled \(T\) bounds its
volume-loss derivative by \((n+1)\langle\gamma^n\rangle\cdot T\).
Integration and vanishing of volume on the big-cone boundary give one
uniform \(r_0>0\) with \(\alpha-r_0T\) big for every allowed model.
Choose \(a\) so \(aT\pm H\) are nef. Thus \(\lambda\alpha\geq\pm H\) for
a uniform \(\lambda\).

Homogeneity and order sandwich \(\langle(\alpha+tH)^n\rangle\) between
\((1\pm\lambda|t|)^n\langle\alpha^n\rangle\). Write \(H=T_1-T_2\) with
fixed nef Cartier classes and pair separately with them. The upper class
\(\gamma\) bounds their pairings. We get
\[|[\langle(\alpha+tH)^n\rangle-\langle\alpha^n\rangle]\cdot H|
\leq K_c|t|,\] uniformly in \(D\) and its model. After the constant
correction \(D+tH-|t|\|h\|F\), the cost is at most \(B'=B+2V\|h\|\).
Integrate the full model derivative and use (9.1) along the segment to
obtain \[|\mathscr E(D+tH)-\mathscr E(D)-t\!\int h\,d\mu_D|
\leq K_ct^2/2+|t|\|h\|(B+B')/c.\tag{9.2}\] The permitted small \(t\)
interval may depend on \(c\). First fix large \(c\), then let either
sign of \(t\) tend to zero, then let \(c\to\infty\).

For normalized genuine \(u\in\mathcal E^1\), direct its nonpositive
rational model majorants \(f\geq u\). Their energies are at least
\(E_\theta(u)\), and their measures have raw mass \(V\). For fixed
\(t\), \(P_\theta(f+th)\) decreases all-point to \(P_\theta(u+th)\): the
genuine limit is below \(u+th\), and every competitor for that obstacle
lies below each member. A genuine lower bound \(u-|t|\|h\|\) prevents
generic collapse. Its energy is finite and the energy decreases to it.
Pass (9.2) through a weak measure cluster point and these energy limits.
Dividing by \(t\), taking both signs and then \(c\to\infty\), identifies
its continuous model integrals with the two-sided outer derivative.
Model density and fixed mass make the cluster point unique. Define that
measure as \(M_\theta(u)\). It is an actual positive Radon measure on
\(K\), before any target is selected, and
\[M_\theta(u)(K)=\operatorname{vol}(L),\qquad
\left.\frac{d}{dt}\right|_0E_\theta(P_\theta(u+th))
=\int h\,dM_\theta(u),\quad h\in C(K).\tag{9.3}\] For continuous \(h\),
uniform model approximation changes a quotient by at most
\(V\|h-h_j\|\), with either sign; thus no original continuous observable
is lost.

The same proof on bounded absolute obstacle sets gives the original
obstacle measure and its uniform-obstacle continuity. To check agreement
at \(u=P_\theta(g)\), take uniformly convergent model upper
approximations \(g_i\) and cofinal model majorants \(f_j\) of \(u\).
Their minima are majorants of \(u\), and their energy gap from \(g_i\)
is at most \(\mathscr E(g_i)-E_\theta(u)\to0\). Using both signs in
(9.2) and obstacle order gives \[\left|\int h(d\mu_f-d\mu_g)\right|
\leq\frac{\mathscr E(g)-\mathscr E(f)}{|t|}
+K_c|t|+\epsilon_c,\quad f\leq g,\quad\epsilon_c\to0.\] Hence ordered
vanishing energy gaps identify the two limiting measures. This proves
agreement with (7.3) for model envelopes; it is not an unrelated
supergradient assignment.

The ordered gap estimate holds at genuine bases by (9.3). It proves weak
measure continuity for decreasing genuine finite-energy nets, using
Section 8 energy continuity. If \(u_i\to u\) divisorially and
\(E_\theta(u_i)\to E_\theta(u)\), the genuine maxima \(\max(u_i,u)\)
converge in energy by upper semicontinuity and order. Apply the ordered
estimate to both endpoints below this maximum. Thus
\(M_\theta(u_i)\to M_\theta(u)\) against all original continuous tests.
Constant shifts leave the operator unchanged.

\section{10. What this construction still has to prove for the
equation}\label{what-this-construction-still-has-to-prove-for-the-equation}

The intrinsic candidate target condition is, for positive Radon \(\nu\)
of raw mass \(V\), \[\nu(\{V_\theta=-\infty\})=0,\qquad
D_\nu=\sup_{\substack{u\in\mathcal E^1(\theta)\\u(v_0)=0}}
\left[E_\theta(u)+\int(V_\theta-u)\,d\nu\right]<\infty.\tag{10.1}\] The
integral is taken off the null reference pole set, and is nonnegative.
This condition mentions the entire actual space (8.2), not a selected
inverse or target-supported skeleton. Mere positivity and correct mass
do not prove (10.1).

To complete P16.5, one must prove, for the independently constructed
operator (9.3), each of the following original-carrier assertions. These
are not conclusions of Sections 3--9 alone:

\begin{enumerate}
\def\labelenumi{\arabic{enumi}.}
\tightlist
\item
  Every output has the reference pole convention and integrates every
  relative genuine \(\mathcal E^1\) test, with finite intrinsic dual
  energy.
\item
  The condition (10.1) is equivalent to integrability of every relative
  genuine finite-energy test; it yields a coercive objective, rather
  than only formal scalar duality.
\item
  The objective has an actual attained maximizer in (8.2) for every such
  arbitrary-support target. Compact energy sublevels alone do not show
  continuity of its unbounded integrals.
\item
  The exact entire supporting plane
  \[E_\theta(v)-E_\theta(u)\leq\int(v-u)\,dM_\theta(u)\] must hold with
  finite relative integrals. Outer derivatives in original continuous
  directions do not by themselves prove this inequality in an unbounded
  genuine direction.
\item
  Comparison on \(\{u<v\}\) and all-solution uniqueness modulo constants
  must be established. Equality of finite model derivatives, metric
  separation, or uniqueness only of variational maximizers is
  insufficient.
\end{enumerate}

The historical chain proposes specific arguments for these points. The
next chapter records their load-bearing statements and the limit orders
that a complete reconstruction must justify. No historical conclusion or
previous review is substituted for those proofs here.

\section{11. A genuine finite-model equation
deduction}\label{a-genuine-finite-model-equation-deduction}

The following conclusion is a proved conditional deduction from Sections
2--9. Its conditionality includes the relative multiplier-ideal
generation and the external algebraic inputs in Section 1. It is not the
whole-support equation requested in original P16.5.

Fix one smooth snc equivariant model with central fiber
\(F=\sum b_iE_i\). Let \(x_i\) be its normalized original-carrier
central valuations. For an arbitrary genuine \(u\), a model obstacle
\(f_D\) with values \(d_i/b_i\) satisfies \[u\leq f_D\text{ on }K
\quad\Longleftrightarrow\quad
u(x_i)\leq d_i/b_i\text{ for all }i.\tag{11.1}\] Here is the all-point
proof. For \(x\in K\), take its Gauss semivaluation \(w\) centered at a
point of the model. Locally \(t\) is a unit times \(\prod z_i^{b_i}\),
with \(z_i\) the central components through that center. Set
\(\lambda_i=b_iw(z_i)\); these are finite nonnegative numbers and sum to
one because \(w(t)=1\). A pulled regular section factors by the
corresponding central prime orders, and its remaining local factor has
nonnegative semivaluation. Hence its normalized logarithm at \(x\) is at
most the weighted sum of its values at the \(x_i\). A weighted section
maximum has the same inequality because the maximum of sums is at most
the sum of maxima. Pass through the original defining FS net; there are
only finitely many vertices and their genuine limiting values are
finite. This gives \[u(x)\leq\sum_i\lambda_i u(x_i).\] The model
function is exactly \(\sum_i\lambda_i f_D(x_i)\), proving (11.1). No
assumption of full support for \(x\) enters this argument.

The model measures have contact orthogonality. Indeed the
continuous-obstacle energy derivative and envelope-fiber descent give
\(\mu_f(\{P_\theta(f)<f\})=0\). To see it directly, choose a nonnegative
continuous test \(h\) supported where the nonnegative lsc gap
\(f-P_\theta(f)\) is larger than a positive number. For sufficiently
small \(t>0\), \(P_\theta(f-th)=P_\theta(f)\). Differentiation gives
\(\int h\,d\mu_f=0\), and Radon regularity, then a countable union over
positive rational gaps, proves the assertion.

Let \(S\) be a nonempty subset of these vertices, \(p_i>0\) for
\(i\in S\), and \(\sum_{i\in S}p_i=V\). Choose the fixed actual ample
growth constants from Section 2 and let
\(T=\max_{i\in S}T_\Omega(x_i)\). For \(a\in\mathbb R^S\), assign
obstacle values \(a_i\) on \(S\) and \(\max_i a_i+T+1\) at every extra
vertex. By the growth bound, any genuine function satisfying the
constraints on \(S\) has generic value at most
\(\min_i(a_i+T_\Omega(x_i))\) and hence is strictly below every extra
obstacle coordinate. Consequently its envelope \(w_a\) imposes exactly
the desired constraints on \(S\) and its measure gives every extra
vertex zero mass.

The genuine real-valued objective
\[G(a)=E_\theta(w_a)-\sum_{i\in S}p_i a_i\] is continuous, because the
obstacle depends continuously in sup norm on finitely many coordinates.
It is invariant under simultaneous constants. Normalize
\(\max_i a_i=0\). Minimality of \(V_\theta\) and the constant energy law
give \(E_\theta(w_a)\leq V\,w_a(v_0)\leq V\min_i a_i+VT\). Thus
\[G(a)\leq VT-p_{\min}\operatorname{osc}(a).\tag{11.2}\] Its upper
levels in the normalized finite-dimensional parameter space are compact,
so a maximum \(a^*\) is attained.

For differentiation, replace the extra obstacle coordinates near any
fixed \(a\) by one sufficiently large fixed constant. Their constraints
are still inactive. The actual obstacle is now affine in the desired
coordinates, and its two-sided derivatives give
\[\frac{\partial}{\partial a_i}E_\theta(w_a)=\mu_a(\{x_i\}).\] At the
maximum all derivatives of \(G\) vanish. Therefore the operator
constructed before this target choice satisfies
\[M_\theta(w_{a^*})=\sum_{i\in S}p_i\delta_{x_i}.\tag{11.3}\] The
solution is an actual original-carrier envelope with
\(V_\theta-C\leq w_{a^*}\leq V_\theta+C\), hence lies in (8.2). In
addition it attains the entire-cone variational supremum for this finite
target: for any genuine \(u\), take \(a_i=u(x_i)\), use \(u\leq w_a\)
and energy order, then maximize \(G\). Contact at every positive target
atom identifies the objective at the maximizer. This proof does not
establish uniqueness among arbitrary entire-space solutions, or
arbitrary-support existence.

\section{12. Preserved counterexamples and normalization
checks}\label{preserved-counterexamples-and-normalization-checks}

\subsection{12.1. Abstract twice-decreasing closure is
false}\label{abstract-twice-decreasing-closure-is-false}

Let \(K_0=\{0\}\cup\{1/r:r\geq1\}\). For \(i,j\geq1\), let \(f_{ij}\)
have its sole pole at \(1/i\), value \(-j\) at the other first \(i\)
isolated points, and value zero at \(0\) and every later point. Each
\(f_{ij}+c\) is continuous into \(\mathbb R\cup\{-\infty\}\). As \(j\)
increases, the functions decrease to \(g_i\), equal to \(-\infty\) at
the first \(i\) points and zero elsewhere. The \(g_i\) decrease to
\(g\), with \(g(0)=0\) and \(g(1/r)=-\infty\) for every \(r\).

Any two members of a decreasing net of the available \(f_{ij}+c\) must
admit a common later member below both. A function with one pole cannot
lie below two functions with different poles. Hence that entire net must
have one fixed pole \(1/i\). Its value at zero equals its value at every
later point \(1/r\), \(r>i\), and remains so in the limit. It cannot
have endpoint \(g\). Thus taking the abstract decreasing-net closure
twice enlarges this family. This is a complete mathematical
counterexample; finite checks in the rerun script are only illustrations
of its pole obstruction. It is not a counterexample to closure for
actual algebraic FS families, to existence in a big class, or to the
equation.

\subsection{12.2. A real FS approximation can lose all finite
energy}\label{a-real-fs-approximation-can-lose-all-finite-energy}

Already in an ample granted class, let \(X=\mathbb P^1\) and
\(L=\mathcal O(1)\), with homogeneous coordinates \(z_0,z_1\). Put
\(\phi=\log|z_1|\) in the canonical metric and \(f_k=\max(\phi,-k)\) for
positive rational \(k\). Appending the two global generators at weight
\(-k\) shows \(f_k\) is a genuine finite FS function on the actual
\(L\).

In degree \(m\), the monomial \(z_1^az_0^{m-a}\) has intrinsic obstacle
weight \[W_m(f_k;z_1^az_0^{m-a})=-k(m-a),\qquad 0\leq a\leq m.\] For
proof, on the original valuation ray
\(v_t=t\operatorname{ord}_{\{z_1=0\}}\), its logarithm is \(-at\) and
the obstacle is \(-\min(t,k)\), so the minimum of their bracket is
attained at \(t=k\) and equals \(-k(m-a)\). Away from that branch the
obstacle is zero and canonical section norms are at most one, giving
nonnegative brackets. At a section pole the bracket is \(+\infty\) as
stipulated. For a general section, its order at the point determines the
same filtration jump; the high-weight spaces are precisely the subspaces
with the corresponding order of vanishing.

Therefore, ignoring the vanishing floor error,
\[e_m(f_k)=-\frac{k}{m^2}\sum_{a=0}^m(m-a)
=-\frac{k(m+1)}{2m}\longrightarrow-k/2.\] The genuine \(f_k\) decrease
all-point to the actual one-section \(\phi\), so Section 8 or ample G0
gives \(E(\phi)=-\infty\). This proves a bad finite-energy approximation
can occur in the granted ample setting itself. It does not imply
nonexistence of finite-energy solutions for a target, and it does not
invalidate the relative-floor repair. It prevents treating arbitrary
uncapped fixed-system FS functions as energy-finite model approximants.

There is also a direct granted-ample curve calculation, so this
counterexample does not depend on the conditional big construction.
Along the one nonconstant valuation ray the slopes of \(f_k\) are \(-1\)
before \(k\) and zero after \(k\). The finite model curve formula gives
\(dd^cf_k=\delta_{v_k}-\delta_{v_0}\), with the canonical trivial
curvature \(\delta_{v_0}\). Thus \(M(f_k)=\delta_{v_k}\) and the raw
degree-one energy is
\[E(f_k)=\frac12\int f_k\,(\delta_{v_0}+\delta_{v_k})=-k/2.\] The ample
decreasing-energy theorem then gives \(E(\phi)=-\infty\) independently
of Sections 3--9.

\subsection{12.3. Raw mass in a big nonnef class is not ordinary top
intersection}\label{raw-mass-in-a-big-nonnef-class-is-not-ordinary-top-intersection}

Let \(X=\operatorname{Bl}_p\mathbb P^2\), \(H\) the pulled hyperplane
and \(E\) the exceptional divisor. For the actual integral \(L=2H+E\),
\[L\cdot E=-1,\qquad L^2=3.\] For \(1\leq j\leq m\),
\((2mH+jE)|_E=\mathcal O_{\mathbb P^1}(-j)\). Its restriction sequence
identifies the section spaces before and after multiplication by the
\(E\) section. Thus \[H^0(X,mL)=s_E^mH^0(X,2mH),\qquad
h^0(X,mL)=\binom{2m+2}{2},\qquad \operatorname{vol}(L)=4.\] This proves
mass \(L^2=3\) would be wrong for this big class. The section-volume
mass in Section 7 is \(4\). The simple integer outputs in the rerun
script check the arithmetic, while the restriction argument proves the
asymptotics.

\section{13. The separate zero numerical class and a zero-volume
obstruction}\label{the-separate-zero-numerical-class-and-a-zero-volume-obstruction}

Retain the same field and smooth projective carrier, and now take any
actual rational bundle with \(c_1(L)=0\), including a nontorsion element
of \(\operatorname{Pic}^0(X)\). Its genuine cone is numerical and uses
varying actual bundles; it is not defined by fixed-degree sections of
this \(L\).

Constants belong to \(\operatorname{PSH}(0)\) directly: the actual
bundles \(L+\epsilon A\) are ample, and weight-zero generators represent
zero as \(\epsilon\downarrow0\). Conversely, let
\(u\in\operatorname{PSH}(0)\) be generically normalized to zero. For
every positive integer \(m\), \(mu\) is genuine in the zero numerical
class and embeds as the same normalized function into a fixed actual
ample \(\Omega\) by Section 2. At any fixed divisorial \(z\),
\[-T_\Omega(z)\leq m u(z)\leq0\quad\hbox{for all }m.\] Hence \(u(z)=0\).
G0 all-point determination applied in this actual ample cone makes
\(u=0\) everywhere. Thus \(\operatorname{PSH}(0)=\mathbb R\).

On this cone define the raw energy and genuine top operator by
\(E_0(c)=0\), \(M_0(c)=0\). These agree with positive ample
perturbations on the actual bundles \(L+\epsilon A\): the curvature
measures of the constant metric are \(\epsilon^n(A^n)\delta_{v_0}\),
converging to zero, and the raw constant energies are
\(c\epsilon^n(A^n)\), also converging to zero. The actual finite-energy
cone is precisely the constants. For a continuous obstacle its genuine
envelope is the constant \(\min_K f\), whose energy is still zero. All
continuous outer energy derivatives and mass comparison expressions
vanish. Domination is asserted here only on the one-point normalized
cone: an unnormalized zero-mass domination statement would fail for two
different constants.

A positive Radon target of mass zero is necessarily zero. Its relative
dual energy is zero and the equation is the bijection
\[\mathcal E^1(0)/\mathbb R\longrightarrow\{0\}.\] All solutions are
actual constants and are unique modulo constants. No probability
normalization or division by zero occurs. The raw rooftop distance
vanishes on unnormalized constants, so it is not a separating metric
there; the normalized quotient is one point.

The restriction \(\theta=0\) matters. On
\(X=\mathbb P^1\times\mathbb P^1\), the actual nef nonample bundle
\(L=p_1^*\mathcal O(1)\) has zero top volume, while its genuine cone has
nonconstant bounded FS functions pulled from the first factor. For
example pull back the preceding \(f_1=\max(\log|z_1|,-1)\). It is zero
at the generic valuation and \(-1\) at the valuation with first-factor
order one, so it is not constant.

The positive model metric on this \(L\) has top mass \(L^2=0\). Its
positive genuine top measure is therefore zero. Equivalently, perturb by
the actual ample \(\epsilon(p_1^*\mathcal O(1)+p_2^*\mathcal O(1))\);
the raw masses are \(2\epsilon(1+\epsilon)\to0\), and positivity forces
every weak limit to be zero. Both the constant and the displayed
nonconstant bounded potential solve the zero top equation, and their
difference is not constant. Thus zero top mass cannot give the
positive-volume uniqueness theorem for all zero-volume nef classes. This
is an actual zero-volume uniqueness obstruction, not an assertion that
no alternative lower-degree equation theory can be constructed.

\section{14. Load-bearing assertions in the proposed entire equation
route}\label{load-bearing-assertions-in-the-proposed-entire-equation-route}

The supplied historical chain is preserved unchanged. The following
audit examines its intended actual moving systems and explains the exact
unproved work in this reconstruction. A historical label does not supply
a premise.

For a model obstacle \(f_D\leq0\) with generically normalized envelope
\(u_D\), choose \(c\) sufficiently large for every lattice summand in
(7.1) to have positive degree. Resolve the base ideal of the
\textbf{complete global} system \(|m\alpha_D(c)|\). Its moving part
divided by \(m\) is a globally generated nef divisor \(C_m(D)\) on the
total space. Independently resolve the base ideal of the complete
original \(|mL|\) and write the actual bundle identity
\[p_m^*L=A_m+T_m,\] with \(A_m\) globally generated nef and \(T_m\)
effective. Every lattice summand is generated near every punctured
fiber, so that generic base ideal is the original one; after removing
it, the remaining central construction should give a bounded
semipositive model potential \(\psi_{m,D}\) on the \textbf{same} actual
\(A_m\), with \[\phi_{m,D}=q_{T_m}+\psi_{m,D},\qquad
\phi_{m,0}=q_{T_m}.\] The divisor logarithm may have poles on the
original carrier, and the \(\phi_{m,D}\) may have infinite negative
original energy, as Section 12.2 shows. They may be used as section
filtrations or original inequalities, not as unproved finite-energy
approximants.

Ordinary Fujita approximation predicts \[A_m^n\to V,\qquad
C_m(D)^{n+1}\to\operatorname{vol}(\alpha_D(c)),\qquad
e_m(\psi_{m,D})=\frac{C_m(D)^{n+1}}{n+1}-cA_m^n\to E_\theta(u_D).\tag{14.1}\]
A further positive-product assertion is needed, not just top-volume
convergence: \[C_m(D)^n\to\langle\alpha_D(c)^n\rangle
\quad\hbox{against every fixed Cartier divisor.}\tag{14.2}\] The
historical argument for (14.2) uses the factorial complete-system order
and the least-upper-bound definition of the positive product. Its
admissible competitors must be actual nef rational bundles with an
actual effective remainder. Approximating such a competitor by an ample
one and using its generated subseries is a promising way to identify the
supremum. Top intersection convergence alone does not prove (14.2), and
order may be paired directly only with nef test divisors. A complete
standalone reconstruction of (14.1)--(14.2), including the actual
compactified moving-model identifications, has not been supplied by the
proved sections of this paper.

The two crucial energy transfers have different scopes. First, a fixed
original continuous majorant \(f\) of a physical moving-system potential
should satisfy the missing-rank trace inequality
\[\mathscr E(f)\geq e_m(\psi)+(\min f)(V-A_m^n).\tag{14.3}\] Here the
original section space receives the injected auxiliary moving-series
space, and every missing original section has allowable weight at least
\(m\min f\). The order of limits is auxiliary section degree first for
\textbf{fixed} \(m\), then \(m\to\infty\) for \textbf{fixed} \(f\), and
finally infimum over all original majorants. This can transfer a lower
energy bound without claiming that uncapped original FS systems have
finite energy.

Second, unbounded original tests need a lower-semicontinuous bridge. Let
\(\phi_{\ell,0}\) be the unweighted complete original \(L\)-section FS
maximum and \(Z_\ell\) its zero base locus. A candidate test is
\[g_\ell(v)=\max(\phi_{\ell,0}-v,0)\quad\text{off }Z_\ell,
\qquad g_\ell(v)=0\quad\text{on }Z_\ell.\] It is nonnegative lsc, but
its increasing limit equals \(V_\theta-v\) only off an output-null
discrepancy set. The requisite discrepancy assertion must be proved in
an actual ample embedding, through the complete negligible-set theorem
and its field/domain hypotheses. Nullness of each individual algebraic
base locus does not imply nullness of this discrepancy set. Likewise
absolute integrability of \(V_\theta\) may fail, so subtraction must
remain relative.

With (14.1)--(14.3), the candidate moment and contraction bounds are
\[\int(V_\theta-v)\,dM_\theta(u)\leq B(v)+nB(u),\tag{14.4}\]
\[B((1-r)V_\theta+ru)\leq2^{n+1}r^2B(u),\qquad B(u)=-E_\theta(u).\tag{14.5}\]
At the bounded auxiliary nef model level, (14.4) follows from a support
plane and the finite mixed-energy identity, and (14.5) from the positive
finite mixed intersections, with the auxiliary generic-value error kept
until it tends to zero. Their original unbounded versions require the
two distinct transfers just described. They are not obtained by
replacing the ample reference zero by a singular \(V_\theta\) in an
ample theorem.

For fixed normalized \textbf{model} endpoints, a further proposed
assertion is the exact weighted moving-series limit
\[\int(\psi_{m,A}-\psi_{m,B})\,d\mu_{m,C}
\longrightarrow\int(u_A-u_B)\,dM_\theta(u_C).\tag{14.6}\] The limit
operator at a model obstacle has finite divisorial support. Actual
section maxima approach its envelope from below at those divisors; if
that approximation is proved, it and usc give continuity of the envelope
at each such point. Uniform boundedness of the auxiliary differences
then lets one localize around the finitely many atoms and control the
complementary mass. This is the reason (14.6) is plausible for model
endpoints. It does not justify pairing a moving unbounded Borel test
with an arbitrary weakly convergent measure.

Actual ample Hölder and strict-concavity estimates can be applied to
these bounded potentials only after perturbing the \textbf{actual}
\(A_m\) by an ample \(\epsilon H_m\). For each fixed \(m\), every
relevant intersection is a finite polynomial in \(\epsilon\). The
correct order is \(\epsilon\downarrow0\) first, then moving-system
\(m\to\infty\). The published constants must remain dimension-only when
these bundles and varieties vary. The rendered primary proof of BJ Lemma
7.30 includes the energy telescoping display with a lower index rendered
as \(1\); the algebraic identity needed here is
\[P^n-Q^n=\sum_{j=0}^{n-1}P^jQ^{n-1-j}(P-Q).\] In particular it cannot
have an empty sum for \(n=1\). The screenshot request failed, so the
recorded observation is about the rendered transcript; no assertion
about the original PDF typography is made. Any application must use the
corrected identity and its full Cauchy--Schwarz proof, not rely on a
malformed formula.

The proposed all-target attainment proof has a sound \emph{conditional}
final step worth retaining. Suppose model maximizing \(u_j\) have
uniformly bounded costs, are model \(I\)-Cauchy, converge divisorially
to an actual genuine \(u\), and satisfy \(E_\theta(u_j)\to e_*\). Weak
upper semicontinuity gives only \(E_\theta(u)\geq e_*\). A fixed model
base has the finite-atomic global plane (8.3), with contact replacing
its obstacle by its envelope. For fixed \(j\), divisorial convergence
permits passing the integral against that one finite atomic measure. If
the model four-point estimate gives
\[\left|\int(u-u_j)\,dM_\theta(u_j)\right|\leq C\delta_j^q,\qquad\delta_j\to0,\]
then that plane forces \(E_\theta(u)\leq E_\theta(u_j)+C\delta_j^q\) and
hence rules out the energy jump. This is a valid conditional deduction.
The actual moving-series estimates, strong fixed-target relative
integration, and their premise boundaries remain to be proved before it
is an attainment theorem.

For all-solution uniqueness one must additionally pass fixed-model
planes to arbitrary bases. The proposed capped tests
\[F_{m,K}=\max(\phi_{m,f},-K),\quad
G_{m,K}=\max(\phi_{m,0},-K)\] are original continuous functions; their
envelope difference must approximate the unbounded relative test in
absolute \(L^1\), uniformly over the specified bounded-cost model
outputs. Outside a fixed closed polar collar the caps are inactive and
the difference of envelopes equals \(F_{m,K}-G_{m,K}\). The proper order
is collar cutoff \(K\) first, section degree \(m\) second, weak-measure
index last. It is a reasonable route, but the uniform absolute errors
must be established component by component before deriving the whole
supporting plane or comparison. Scalar signed integral convergence alone
does not give that absolute \(L^1\) estimate.

Finally the roof/max energy inequality and the original continuous outer
derivative would give a local winning measure inequality on
\textbf{open} winning sets by variation. Strict sets for general usc
potentials need not be open. Passing to them requires the preceding
absolute continuous-test approximation and bounded Lipschitz cutoffs,
with the cutoff kept fixed through the weak limit and sent to zero only
afterward. This final comparison deduction is conditional on the missing
weighted-identification theorem, not a direct restriction of a weak
measure limit to a Borel set.

\section{15. Exact retained result and exact unfinished
result}\label{exact-retained-result-and-exact-unfinished-result}

The new result of this reconstruction is the actual reference-floor
mechanism in Sections 3--6, with its explicit ordinary relative
multiplier-ideal premise and a complete all-point directed-net proof.
The compactness use of divisorial nonnegligibility has been
independently traced through its complete primary ample proof and G0
regularization. Sections 7--9 give a detailed conditional reconstruction
of the section-volume entire outer operator on the original genuine
cone, with correct raw volume mass and every original continuous
direction. Section 11 gives a conditional finite-model equation and
actual variational attainment. Section 13 proves the separate zero
numerical-class cone and zero-mass quotient equation using only actual
ample G0 and elementary genuine constructions.

The arbitrary-dimensional, arbitrary-support big equation theorem
required by original P16.5 is \textbf{not proved here}. No unconditional
original-scope canonical result is asserted. To turn the conditional
foundation into an independently checked theorem, the relative
arbitrary-field generation/local-vanishing application in Section 3
needs its complete primary proof or a complete independent replacement.
To turn the operator construction into the equation, (14.1)--(14.6),
output/target pole-nullness and strong relative integration, the full
unbounded supporting plane, actual all-target attainment, and the
absolute comparison passage require complete standalone proofs with
their exact source applications. The supplied historical chain contains
proposed arguments for them; it has not been promoted into a grant.

These gaps are statements of unfinished proof work, not counterexamples
to a general big-class equation theorem. The abstract closure
counterexample, the actual bad FS sequence, and the zero-volume
uniqueness obstruction are each confined to their proved domains.
General nonbig nonzero pseudoeffective classes have no equation
bijection established by this paper. Every zero-mass statement retains
zero mass; no zero measure has been normalized to a probability.

\section{References and source
boundary}\label{references-and-source-boundary}

\begin{enumerate}
\def\labelenumi{\arabic{enumi}.}
\tightlist
\item
  S. Boucksom and M. Jonsson, \emph{Global pluripotential theory over a
  trivially valued field}, arXiv:1801.08229v4, March 2022. The actual
  ample supporting proofs and the precisely recorded algebraic FS/model
  units are used. The abstract Lemma 4.6 and general big consequences
  depending on it are excluded.
\item
  S. Boucksom and M. Jonsson, \emph{Addendum to the article Global
  pluripotential theory over a trivially valued field},
  arXiv:2206.07183v2, September 2023. The one model-obstacle
  section/multiplier-ideal net is reconstructed; its general
  decreasing-closure and compactness consequences are not imported.
\item
  S. Boucksom, C. Favre, and M. Jonsson, \emph{Differentiability of
  volumes of divisors}, arXiv:math/0608260v2, September 2007. Ordinary
  algebraic positive products and big volume differentiation are
  explicit inputs, in the algebraically closed characteristic-zero
  domain.
\item
  J.-P. Demailly, L. Ein, and R. Lazarsfeld, \emph{A subadditivity
  property of multiplier ideals}, arXiv:math/0002035v2, April 2000. Its
  complex absolute domain is recorded; it is not silently used as a
  relative arbitrary-field generation theorem.
\item
  The exact supplied G0 actual-ample body, reproduced in the preserved
  evidence. Historical documentary bindings identify this body but are
  not mathematical proofs of nonample extensions.
\end{enumerate}

The evidence ledger lists actual retrieved coordinates, including failed
source requests, and does not claim whole-source or recursively complete
proof reading. Compilation, checksums, and finite examples have no
bearing on acceptance of the mathematical claims.

\end{document}
